\chapter{Vectors, Spaces and Hilbert Space}\label{ch:hilbert}

\begin{etymology}
\textbf{Algebra} comes from the Arabic \emph{al-jabr}, ``restoration''. It is a
word from the title of al-Khwarizmi's treatise on equations, written in Baghdad
around 820: \emph{al-jabr} was the operation of moving a subtracted term to the
other side of an equation, restoring the balance.

\textbf{Vector} is Latin for ``carrier'', from \emph{vehere}, ``to carry''.
Something that carries you from one point to another.

\textbf{Space} comes from \emph{spatium}, ``extent'', ``interval''. The room in
which things can exist and move.

\textbf{Phase} comes from the Greek \emph{phásis}, ``appearance''. How something
shows itself at a given moment; the moon has phases. For anything that repeats,
it means where in the cycle it is.
\end{etymology}

The state of a qubit is a vector. Before that sentence can mean anything, you need
to know what a vector is, which turns out to be much more general than an arrow,
and what extra structure a space of vectors needs to talk about probabilities.
The answer is a Hilbert space, and getting there took about a thousand years of
algebra.

\section{From equations to structures}

\begin{history}{around 820 \textperiodcentered{} Baghdad}
Muhammad ibn Musa al-Khwarizmi wrote a treatise on solving equations of the
second degree, without symbols, in words: ``a square and ten roots are equal to
thirty-nine''. Today that is $x^2 + 10x = 39$. The answer was less important than
the method. Al-Khwarizmi noticed that the same two operations, moving terms across
(\emph{al-jabr}) and combining like terms (\emph{al-muqabala}), solve every
equation of a given type, whatever the numbers are.
\end{history}

That is the seed of algebraic thinking: the structure of the manipulation
matters more than the particular numbers. Symbols came slowly. François Viète,
around 1591, was the first to use letters systematically for unknown quantities.
René Descartes, in 1637, fixed the convention still used: letters from the end of
the alphabet ($x, y, z$) for unknowns, from the start ($a, b, c$) for known
quantities.

\begin{history}{30 May 1832 \textperiodcentered{} Paris}
Évariste Galois, 20 years old, spent the night before a duel writing a letter that
summarized his work on why equations of degree five have no general solution by
radicals. The answer, he had found, depends on the symmetries of the roots, not
on the roots themselves. To describe those symmetries he used what we now call
groups. He died of his wound the next day. His papers were published in 1846.
\end{history}

After Galois, algebra stopped being about finding numbers and became the study of
structures: a set, some operations, and the rules they obey. Groups have one
operation. Fields, like the rationals, the reals and the complex numbers, have
addition and multiplication with all inverses; Richard Dedekind introduced them
in 1871 and Heinrich Weber gave the abstract definition in 1893. And
\emph{vector spaces} have addition plus multiplication by numbers taken from a
field.

\section{What a vector is}

Vectors first appeared in physics, for quantities with size and direction: force,
velocity. Hermann Grassmann, a schoolteacher in Stettin, built a general calculus
of directed quantities in 1844; almost nobody read it. Josiah Willard Gibbs and
Oliver Heaviside turned vectors into the everyday tool of physics in the 1880s.
Giuseppe Peano, in 1888, gave the axiomatic definition, and with it the key
step: what makes something a vector is not that it is an arrow, but the rules it
obeys.

\begin{concept}{Vector space}
A set of objects (vectors) that can be added to each other and multiplied by
numbers (scalars), where the result is always in the set and the usual rules
hold: addition is commutative and associative, there is a zero vector, every
vector has an opposite, and scaling distributes over addition.

\textbf{What it fixes:} any theorem proved for vector spaces applies at once to
arrows, lists of numbers, polynomials, functions and quantum states.

\textbf{What it costs:} the intuition of arrows in the plane stops being reliable.
You have to check the rules, not the picture.
\end{concept}

Arrows in the plane are vectors. So are lists of $n$ real numbers, polynomials of
degree at most $n$, continuous functions, and lists of $n$ complex numbers, which
form the space $\mathbb{C}^n$.

The rule that is easiest to forget is \emph{closure}: adding two vectors, or
scaling one, must give something inside the set. Points of the unit disk are not
a vector space: $(0.8, 0) + (0.8, 0) = (1.6, 0)$ falls outside. The positive
quadrant, where both coordinates are positive, is closed under addition (two
positive numbers add to a positive number) but not under scaling: multiply
$(1, 2)$ by $-1$ and you leave it. Without closure, none of the theorems of linear
algebra apply.

\section{Complex numbers}

\begin{history}{1545--1572 \textperiodcentered{} Milan and Bologna}
Girolamo Cardano published the solution of the cubic equation in the \emph{Ars
Magna} of 1545. For some equations with three perfectly real solutions, the
formula passes through square roots of negative numbers along the way. Cardano
called these quantities ``as subtle as they are useless''. Rafael Bombelli, in his
\emph{Algebra} of 1572, worked out rules to compute with them, and showed that
they cancel out and leave the real answer. Complex numbers were not invented to
solve $x^2 = -1$. They were forced on mathematicians by equations that had real
solutions.
\end{history}

For two centuries, these ``imaginary'' numbers were used without being
understood. Caspar Wessel in 1799, Jean-Robert Argand in 1806 and Carl Friedrich
Gauss in 1831 gave them a geometric meaning: $z = a + ib$ is a point in the plane,
with $a$ on the horizontal axis and $b$ on the vertical one. Leonhard Euler had
already found, in 1748, the identity that makes the picture work:
\[
  e^{i\theta} = \cos\theta + i\sin\theta .
\]
So every complex number can be written as $z = re^{i\theta}$, with a size $r$ and
an angle $\theta$, the \emph{phase}. Multiplying by $e^{i\theta}$ rotates by
$\theta$; multiplying by $i$ rotates by 90 degrees. The algebra of complex numbers
was built so that multiplying means rotating.

In $\mathbb{R}^n$, scaling a vector can only stretch it or flip its sign. In
$\mathbb{C}^n$, scaling can also rotate it. A real number has two directions,
positive and negative; a complex number has a whole circle of them.

Angles exist in $\mathbb{R}^2$ too, of course, but there they are something you
compute between two vectors. In $\mathbb{C}$ the angle is part of the number
itself, and it travels with every scalar you use.

You can describe oscillations with real numbers only: $A\cos(\omega t + \theta)$
has an amplitude, a frequency and a phase. But shifting a phase or adding two
oscillations then needs trigonometric identities, and the bookkeeping grows fast.
Written as $Ae^{i(\omega t + \theta)}$, an oscillation is a rotating point, and
shifting its phase is one multiplication. Real numbers can do the job; complex
numbers do it with the structure built in. Quantum amplitudes oscillate, and
interference depends on their phases, which is why they are complex.

\section{Hilbert space}

\begin{history}{1904--1929 \textperiodcentered{} Göttingen and Berlin}
Between 1904 and 1910, David Hilbert worked on integral equations, which appear
all over physics. They could not be handled with finite lists of numbers: the
unknowns were functions, which live in spaces of infinite dimension. Hilbert
showed that length, angle and perpendicularity still make sense there, through
infinite sequences whose squares add up to a finite number. In 1929 John von
Neumann, who had worked with Hilbert in Göttingen, gave the abstract definition
and the name: Hilbert space. He built quantum mechanics on it three years later.
\end{history}

A Hilbert space is a vector space over $\mathbb{C}$ with an \emph{inner product}
$\braket{u|v}$, a way of combining two vectors into one number, that satisfies
four conditions:

\begin{enumerate}
\item it is linear in the second argument;
\item swapping the arguments gives the complex conjugate,
$\braket{u|v} = \overline{\braket{v|u}}$;
\item $\braket{v|v} \geq 0$, and it is zero only for the zero vector;
\item the space is \emph{complete}: every Cauchy sequence converges to a point
inside the space.
\end{enumerate}

A Cauchy sequence, named after Augustin-Louis Cauchy, is one whose terms get
closer and closer to each other, like $1 + 1/n$. Completeness guarantees that
the point they are approaching exists in the space. The rationals are not
complete: a sequence of fractions can approach $\sqrt{2}$, which is not a fraction.
In the finite spaces of quantum computing completeness is automatic, but it is
what makes the infinite-dimensional spaces of physics usable.

In $\mathbb{C}^n$ the inner product is
$\braket{u|v} = \sum_i \overline{u_i}\,v_i$. The conjugation on the first vector is
what makes $\braket{v|v} = \sum_i |v_i|^2$ real and non-negative. Without it, the
``length'' of a complex vector could come out negative or complex.

Everything geometric comes out of this one operation:

\begin{center}
\begin{tabular}{lll}
\hline
concept & from the inner product & meaning \\
\hline
norm & $\|v\| = \sqrt{\braket{v|v}}$ & the length of the vector \\
orthogonality & $\braket{u|v} = 0$ & the vectors do not overlap \\
projection & $\dfrac{\braket{u|v}}{\braket{u|u}}\,u$ & how much of $v$ lies along $u$ \\
\hline
\end{tabular}
\end{center}

\section{The state vector}

A quantum state is a vector $\ket{\psi}$ in a Hilbert space with norm 1:
$\braket{\psi|\psi} = 1$. The system can be a particle's position, a photon's
polarization, or a qubit. For one qubit,
\[
  \ket{\psi} = \alpha\ket{0} + \beta\ket{1}, \qquad |\alpha|^2 + |\beta|^2 = 1 .
\]
The state does not store physical values directly. It stores how the system will
answer when you ask it something: each basis vector is a possible answer, and each
coefficient says how much the state is aligned with it.

Why a vector and not something else? Superposition is linear: if two states are
possible, any normalized combination of them is too, and that is exactly what a
vector space guarantees. Probabilities come from inner products,
$P = |\braket{\phi|\psi}|^2$, and the norm makes them add to 1. And operations on
states (gates, time evolution, measurements) are linear maps on the space.

The state vector is not a probability distribution: the amplitudes are complex,
and their phases, which do not show up in $|\alpha|^2$, decide interference. It is
not the physical system either, only its description. And not all of its phase is
physical.

\begin{concept}{Global and relative phase}
Multiplying the whole state by $e^{i\gamma}$ changes no probability, because
$|e^{i\gamma}\braket{\phi|\psi}|^2 = |\braket{\phi|\psi}|^2$. The states
$\ket{\psi}$ and $e^{i\gamma}\ket{\psi}$ are physically the same. But a phase
between components, as in $\ket{0} + e^{i\varphi}\ket{1}$, changes how they
interfere.

\textbf{What it fixes:} it says which part of the description is physics and which
is bookkeeping.

\textbf{What it costs:} the relative phase is invisible to a direct measurement
and decisive for the next operation. Quantum algorithms work by manipulating
exactly that.
\end{concept}
